\chapter{Ideally exact categories}


\section{Preliminary results on the Eilenberg-Moore categories}

We present here some technical results on the Eilenberg-Moore categories. These will play a key role in the next section, where we define and characterize the ideally exact categories through monadic functors.

\begin{mainbox}{}
    \begin{prop}\label{algebras category is exact if the base category is exact}
        Let $T:\mathcal{\mathcal{X}} \to \mathcal{X}$ be a monad preserving regular epimorphisms and kernel pairs. If $\mathcal{X}$ is exact, then the Eilenberg-Moore category $\mathcal{X}^T$ is exact.
    \end{prop}
\end{mainbox}

\begin{subbox}{}
    \begin{proof}
        Assume $\mathcal{X}$ is exact. With the given hypotheses on the monad $T$, we aim to prove that $\mathcal{X}^T$ is exact. We need to establish that it is finitely complete, that each equivalent relation is the kernel pair of some arrow, that regular epis are pullback stable and that each morphism admits a (regular epi, mono) factorization.

        First we show that $\mathcal{X}^T$ is finitely complete. We recall that the forgetful functor $U^T:\mathcal{X}^T \to \mathcal{X}$ creates all the limits that $\mathcal{X}$ has. Since $\mathcal{X}$ is exact, it has all finite limits. Therefore, $\mathcal{X}^T$ has all finite limits.

        We proceed to prove that each equivalent relation is the kernel pair of some arrow. Consider an equivalence relation $R$ in $\mathcal{X}^T$ with $\xi_R$ and $\xi_A$ giving the structure of algebras. Then, under the forgetful functor $U^T$, $R$ is an equivalence relation in $\mathcal{X}$. $\mathcal{X}$ is exact, so $R$ is the kernel pair of some arrow. In particular, by \Cref{equiv rel is kernel pair of the coeq}, we cas suppose that $R$ is the kernel pair of its coequalizer. Let call it $f:A \to B$. Since $T$ preserves kernel pairs, we have that $T(R)$ is the kernel pair of $T(f)$. Morover $T$ preserves regular epis, so $T(f)$ is a regular epi. By \Cref{coeq of the kernel pair}, we can assume that $T(f)$ is the coequalizer of its kernel pair $(T(R),r_1,r_2)$. Since $f\xi_AT(r_1)= fr_1\xi_R = fr_2\xi_R=f\xi_AT(r_2)$, by the universal property of the coequalizer, there exists a unique morphism $\xi_B$ such that $\xi_BT(f) = f\xi_A$. The morphism $\xi_B$ is the algebra structure of $T(B)$. Since $U^T$ creates all the limits that $\mathcal{X}$ has, we can conclude that $R$ is the kernel pair in $\mathcal{X}^T$ of some arrow.
        \[
        % https://tikzcd.yichuanshen.de/#N4Igdg9gJgpgziAXAbVABwnAlgFyxMJZABgBoAmAXVJADcBDAGwFcYkQAlEAX1PU1z5CKchWp0mrdgEEefEBmx4CRACxiaDFm0QgAQnP5KhRMsXFapugCoAKDgEpDCgcuHJR5zZJ0g70p14jQRUUdS8JbXY7PUDxGCgAc3giUAAzACcIAFskAEYaHAgkUUirEDSQGgALGHoodhwAdwha+oQgiqzcxDIQIvyaOGqsNJwkAFoCst8MgH0850ycpD6BxGnh0fGN7yjdefIqkEZ6ACMYRgAFVxMDrETq8c7lnoBmQuLevfKAHV+AB5YOZcGiMLBgXxwCDghovbpIdT9L7TSy+f5AuayeErRBI9YAVh+vjsaScNTqDV0zValI68leSCJyJKxPYGOBBjBEN8UHowwSSwRiA+LLxQxGY0mqJ80Vs8zygQZwtF6yRWyluxO50uN2MoRAGQeT2OaLlhzi3CAA
\begin{tikzcd}
    T(R) \arrow[dd, "\xi_R"] \arrow[rr, "T(r_1)", shift left] \arrow[rr, "T(r_2)"', shift right] &  & T(A) \arrow[dd, "\xi_A"] \arrow[rr, "T(f)", two heads] &  & T(B) \arrow[dd, "\xi_B", dashed] \\
                                                                                                 &  &                                                        &  &                                  \\
    R \arrow[rr, "r_1", shift left] \arrow[rr, "r_2"', shift right]                              &  & A \arrow[rr, "f", two heads]                           &  & B                               
    \end{tikzcd}
    \]

        Consider now a morphism $f:(A,\xi_A)\to (B, \xi_B)$ in $\mathcal{X}^T$, we aim to prove that it admits a (regular epi, mono) factorization. The morphism $f$, under the forgetful functor $U^T$, is a morphism in $\mathcal{X}$, which is exact and in particular regular. Thus $f$ admits a (regular epi, mono) factorization in $\mathcal{X}$. Let it be $f=me$ with $m:Q \to B$ mono and $e: A \to Q$ regular epi. Notice that $A$ and $B$ are already algebras, we want to provide $Q$ and the domain of the parallel arrows with a structure of algebra, so that we can lift the factorization to $\mathcal{X}^T$.
        \[
% https://tikzcd.yichuanshen.de/#N4Igdg9gJgpgziAXAbVABwnAlgFyxMJZAJgBpiBdUkANwEMAbAVxiRAEEQBfU9TXfIRQAWclVqMWbAELdeIDNjwEiAZlKrx9Zq0QgAinL5LBRMgAYtk3SAAqACnYBKIwv7KhyUZerapeh2kXHmMBFRRzMV9rNgBhV0Uwz0ifCR02ABFucRgoAHN4IlAAMwAnCABbJEiQHAgkMjT-EFZqBjoAIxgGAAV3Uz1SrDyACxwQahGYOig2HAB3CCmZhBCQMsqG6jqkAEZtuiwGOcPj6PS9KrbO7r6TcJAh0fG1jarEGp3EfaabYtc3kh1LV6h9zs0ADoQgAeWAA+pxXuV3qIQXtwTYobC4bIkZtEMCvqi-DYHMVgvJAYgAKzbUE1OAjLDFcbfPHvABsdKB1EZzNZu3ZSC5aIJvKZLKQAFpBZTkUhaaKGRLWTKuBQuEA
\begin{tikzcd}
D \arrow[rr, shift right] \arrow[rr, shift left] &  & T(A) \arrow[dd, "\xi_A"] \arrow[rr, "T(f)"]   &                          & T(B) \arrow[dd, "\xi_B"] \\
                                                 &  &                                               &                          &                          \\
C \arrow[rr, shift right] \arrow[rr, shift left] &  & A \arrow[rd, "e"', two heads] \arrow[rr, "f"] &                          & B                        \\
                                                 &  &                                               & Q \arrow[ru, "m"', tail] &                         
\end{tikzcd}
        \]

        In order to do so, we first need to show that $T$ preserves coequalizers in $\mathcal{X}$ of kernel pairs arising from $\mathcal{X}^T$, once the algebra structure is forgotten, and considering only the pair of the whole kernel pair. Let $(\mathrm{Eq},r_1,r_2)$ be a kernel pair in $\mathcal{X}$, arising from a kernel pair in $\mathcal{X}^T$ of a morphism, and let $p$ be the coequalizer in $\mathcal{X}$ of $(r_1,r_2)$. By \Cref{kernel pair of the coeq}, we can affirm that $(r_1,r_2)$ is the kernel pair of $p$. Given that $(\mathrm{Eq},r_1,r_2)$ is a kernel pair arising from $\mathcal{X}^T$, it follows that $(T(\mathrm{Eq}),T(r_1),T(r_2))$ is the kernel pair of some arrow. Since $T$ preserves regular epi and kernel pair we can assert that $T(p)$ is a regular epi and, by \Cref{coeq of the kernel pair}, it is the coequalizer of the kernel pair $(T(\mathrm{Eq}),T(r_1),T(r_2))$, as claimed.
        \[
        % https://tikzcd.yichuanshen.de/#N4Igdg9gJgpgziAXAbVABwnAlgFyxMJZABgBpiBdUkANwEMAbAVxiRABUAKAHW4Fs6OABYAnPsACiARwC+AShAzS6TLnyEUZAExVajFm14DhYybMXKQGbHgJEtpHdXrNWiEAEELKm+qIAWR10XA3cAIW8rVVsNZAdKZ303Dk4PBSUfNTsUQIS9VzYuMPTdGCgAc3giUAAzEQg+JDIQHAgkQJA4ISwanCQARkSC9y4RAH0tBWoGOgAjGAYABWi-dxEscqE+jJA6hqbqVvbqLp6+xABaQfzQlPH+9Ms9xsRro8QHTu7egaHb8a0IGmcwWy182RA6022ye9RebzaHxO33OVz+yXukWeBxaiOuDCwYGScAgBKgQJuyV4AA8sGMjIJROJpDIsXDjrikJ8QlTuLSxl4dtjEABWQ6IgDM6MMfLpESF7MQHXeYsphU4aCmICEMDo5PcOAA7hAdXqEAr9kjOYgpWr3GgKab9S1jU7zRQZEA
        \begin{tikzcd}
        T(\mathrm{Eq}) \arrow[rr, "T(r_2)"', shift right] \arrow[rr, "T(r_1)", shift left] \arrow[dd, "\xi_\mathrm{Eq}"] &  & T(A) \arrow[dd, "\xi_A"] \arrow[rr, "T(p)", two heads] &  & T(B) \arrow[dd, "\xi_B"] \\
                                                                                                                         &  &                                                        &  &                          \\
        \mathrm{Eq} \arrow[rr, "r_2"', shift right] \arrow[rr, "r_1", shift left]                                        &  & A \arrow[rr, "p", two heads]                           &  & B                       
        \end{tikzcd}
        \]
        
        Picking up where we left off, consider $f$ and its kernel pair $(\mathrm{Eq}(f),r_1,r_2)$. Since $U^T$ creates all the limits that $\mathcal{X}$ has, we can lift the kernel pair $(\mathrm{Eq}(f),r_1,r_2)$ to a kernel pair in $\mathcal{X}^T$. Let $\xi_\mathrm{Eq}$ be the structure of algebra of $\mathrm{Eq}$.

        We know that $U^T$ creates all the colimits that $\mathcal{X}$ has and $T$ preserves. Hence, $U^T$ creates all coequalizers that $\mathcal{X}$ has of kernel pairs arising from $\mathcal{X}^T$, once the algebra structure is forgotten. Observe that $e$ is such a coequalizer, in fact it is a regular epi and by \Cref{f and e share the kernel pair} $f$ and $e$ share the same kernel pair, because of the factorization $f=me$ with $m$ mono, thus $(\mathrm{Eq}(f),r_1,r_2)$ is also the kernel pair of $e$. Hence, by \Cref{coeq of the kernel pair}, $e$ is the coequalizer of $(\mathrm{Eq}(f),r_1,r_2)$. In particular by the creation of such colimit we proved that $Q$ has the structure of algebra, called $\xi_Q$, and that $e$ can be lifted to a coequalizer in $\mathcal{X}^T$.
        \[
% https://tikzcd.yichuanshen.de/#N4Igdg9gJgpgziAXAbVABwnAlgFyxMJZABgBpiBdUkANwEMAbAVxiRABUAKAHW4Fs6OABYAnPsACiARwC+AShAzS6TLnyEUZAExVajFm14DhYybMXKQGbHgJEtpHdXrNWiEAEELKm+qIAWR10XA3cAIW8rVVsNZAdKZ303Dk4PBSUfNTsUQIS9VzYuMPTLayzYgGZSAEZgpMLOAEUSzJiiKoq6gvdGxV0YKABzeCJQADMRCD4kMhAcCCRAkDghLDGcJGrE7pSRAH0tBWoGOgAjGAYABWi-dxEsQaENjJAJqZnqecXqFbWNxAAtFt8qFdntqi1XpNpohgV9EA5lqt1pttqD9loQMczhdrr5siB7o9npY3jC4QsET9kf8gWjkvtqpEyR85pTgQwsGBklA6CsBliQcleAAPLB7IyCUTiaQyZnQ75spCIkLC7hivZeF4sxBLeEANmoQhgdCgbBwAHcIMbTQh6Q0YEcQCdzlcbgSGDAUdqFYhDUrEABWe3uLh8SE6xHwgDsIZArGxrrx5TYRKegptZvclutJqgCB970QsYDVRAmaQYCYDAYnzoWAY5vrjbj00TuPdGkJD3ThZhwdLcdF4oifaQ-pjQ-V4t6Y91n0pA9VDTGEd9UcpZeX7jGfRkQA
\begin{tikzcd}
    T(\mathrm{Eq}(f)) \arrow[rr, "T(r_2)"', shift right] \arrow[rr, "T(r_1)", shift left] \arrow[dd, "\xi_\mathrm{Eq}(f)", dashed] &  & T(A) \arrow[dd, "\xi_A"] \arrow[rd, "T(e)"', pos=0.7, two heads] \arrow[rr, "T(f)"] &                                             & T(B) \arrow[dd, "\xi_B"] \\
                                                                                                                             &  &                                                                           & T(Q) \arrow[ru, "T(m)"', pos=0.3] &                          \\
    \mathrm{Eq}(f) \arrow[rr, "r_2"', shift right] \arrow[rr, "r_1", shift left]                                                &  & A \arrow[rd, "e"', two heads] \arrow[rr, "f", pos=0.25]                             &                                             & B                        \\
                                                                                                                             &  &                                                                           & Q \arrow[from=uu, "\xi_Q", pos=0.75, crossing over, dashed] \arrow[ru, "m"', tail]                    &                         
    \end{tikzcd}
\]

        In order to prove the factorization of $f$ in $\mathcal{X}^T$, it remains to establish that $m$ is a morphism in $\mathcal{X}^T$ and that it is a mono. Notice that $m$ is a mono in $\mathcal{X^T}$, in fact it is a mono in $\mathcal{X}$ and furthermore it is merely a morphism $\mathcal{X}$ that satisfies the algebra condition. We still need to establish that $m$ is a morphism in $\mathcal{X}^T$.
        \begin{align*}
            f=mq &\implies f \xi_A = mq\xi_A \\[1.5ex]
            &\implies \xi_B T(f) = m \xi_Q T(q) \\[1.5ex]
            &\implies \xi_B T(m) T(q) =  m \xi_Q T(q)
        \end{align*}
        We know that $T(q)$ is a regular epi, thus it is a epi. So we can cancel it from both sides and we obtain the condition for $m$ to be a morphism in $\mathcal{X}^T$:
        \[ \xi_B T(m) =  m \xi_Q \]
        We proved that $m$ and $p$ are respectively a mono and a regular epi in $\mathcal{X}^T$. This allows the lifting of the factorization $f=me$ to $\mathcal{X}^T$.

        It remains to prove the last condition of a regular category, that is the pullback stability of regular epi. Let $q: A \to B$ be a regular epi in $\mathcal{X}^T$, ie a coequalizer of some arrows in $\mathcal{X}^T$. In particular, $q$ is a regular epi in $\mathcal{X}$. Since $\mathcal{X}$ is exact, it is pullback stable. Let $\tilde{q}$ the pullback in $\mathcal{X}^T$ of $q$ along some morphism. As just noticed, $\tilde{q}$ is a regular epi in $\mathcal{X}$. Let $(\mathrm{Eq}(\tilde{q}), r_1,r_2)$ be the kernel pair of $\tilde{q}$. By \Cref{coeq of the kernel pair}, $\tilde{q}$ is the coequalizer of such kernel pair. We want to lift $\tilde{Q}$ to a regular epi in $\mathcal{X}^T$. By construction, the pullback considered is a pullback in $\mathcal{X}^T$, thus it comes with algebra structures of $A\times_BC$ and $C$. Let $\xi_{A\times_BC}$ and $\xi_C$ be the algebra structures of $A\times_BC$ and $C$ respectively. Moreover, $\tilde{q}$ is a morphism in $\mathcal{X}^T$. We know that the forgetful functor $U^T$ creates all the limits that $\mathcal{X}$ has, so we can lift the kernel pair $(\mathrm{Eq}(\tilde{q}),r_1,r_2)$ to a kernel pair in $\mathcal{X}^T$. Let $\xi_{\mathrm{Eq}(\tilde{q})}$ be the algebra structure of $\mathrm{Eq}(\tilde{q})$.
        We know that $U^T$ creates all the colimits that $\mathcal{X}$ has and the monad $T$ preserves, and we already proved that $T$ preserves coequalizers of kernel pairs. Thus, since $\tilde{q}$ is such a coequalizer, we can lift it to a coequalizer in $\mathcal{X}^T$.
        \[
% https://tikzcd.yichuanshen.de/#N4Igdg9gJgpgziAXAbVABwnAlgFyxMJZAZgBoAmAXVJADcBDAGwFcYkQBBEAX1PU1z5CKAKwVqdJq3YAhHnxAZseAkTEAGCQxZtEIACoAKGQEp5-ZUKJlNNbdL1GOZ3hcGqUARnF2pukADS5ooCKsLI3raSOuxGAS4KSu7h5KTEWn7sHAA62XgAtvAA+jIAwsFJYUQALGkZMXrlriGWHsipnvUOBoY5eViFcCWlCW5VKLWdvg09IxWhVijqddPdufn0OAAWAE75wACiAI7chrl4jLDAJ6MtyUTLU9HdRuubu-vHp+dYlzDX3BMLgkMCgAHN4ERQAAzHYQfJIAAcNBwECQxGasPh6JRaMQqWe-iMRzMNC2MHoUCQYGYjEYmLhCMQyxAqKQ3kJ7COIDJFKpehwAHcIOTKQgGdjEAA2XFIdQSpkATlliAA7Kt-LkAB5YIpNBRYpky1l49UgUX81nCi0IDXsH5-AE8kCMegAIxgjAACgsPCAdlgwVscMFDUiVcb7JrsjqisA+gVimVuKHGUgCWzEJ4FeGTUhlea+dTaYw7Y4zv1HTdU5KC5nyDnEGQ88zG2IW824FssNCQ4gALTZg1pxDtzOd7u99k0V0e72+4T+wPBmtM2otlldnt9weN9eZzeTvsc2een2tRcBoMhtsq6qNzwczOeFlR+0x3XAN7bPaHE4Vi4rhuFMZywMB-Cgegu1BVd2SfPFkRALcpwHIcYRHR8VUQ5DjwfFlM2NHDp05PQdiKchnVPecL3YK8V0bM1nzwiMaCI1Cy39cieEobggA
\begin{tikzcd}
    & T(K) \arrow[rr, shift left] \arrow[rr, shift right] \arrow[dd] &                                                                                    & T(A) \arrow[rr, "T(q)"] \arrow[dd] &                                     & T(B) \arrow[dd] \\
T(\mathrm{Eq}(\tilde{q})) \arrow[dd, "\xi_{\mathrm{Eq}(\tilde{q})}", dashed] \arrow[rr, crossing over, "T(r_1)", pos=0.25, shift left] \arrow[rr, crossing over, "T(r_2)"', pos=0.25 , shift right] &                                                                & T(A\times_BC) \arrow[ru]  \arrow[rr, crossing over, "T(\tilde{q})", pos=0.25] &                                    & T(C) \arrow[ru] &                 \\
    & K \arrow[rr, shift left] \arrow[rr, shift right]               &                                                                                    & \ A \ \arrow[rr, "q", pos=0.25, two heads]       &                                     & B               \\
\mathrm{Eq}(\tilde{q}) \arrow[rr, "r_2"', shift right] \arrow[rr, "r_1", shift left]                                        &                                                                & A\times_BC \arrow[from=uu, crossing over, pos=0.25,  "\xi_{A\times_BC}"] \arrow[ru, shift right] \arrow[rr, "\tilde{q}"', two heads]                          &                                    & C \arrow[from=uu, crossing over, "\xi_C", pos=0.25] \arrow[ru]                        &                
\end{tikzcd}
\]

    \end{proof}
\end{subbox}

\begin{mainbox}{}
    \begin{prop}\label{algebras category is proto if the base category is proto}
        Let $T:\mathcal{\mathcal{X}} \to \mathcal{X}$ be a monad. If $\mathcal{X}$ is protomodular, then the Eilenberg-Moore category $\mathcal{X}^T$ is protomodular.
    \end{prop}
\end{mainbox}

\begin{subbox}{}
    \begin{proof}
        Since the forgetful functor $\mathcal{X}^T \to \mathcal{X}$ creates all the limits that $\mathcal{X}$ has, we can conclude that $\mathcal{X}^T$ is finitely complete.

        Let now $f:(X,\xi) \to (Y,\xi')$ be a morphism in $\mathcal{X}^T$. The following diagram of funtors is commutative. This essentially follows from the fact that the forgetful functor $\mathcal{X}^T \to \mathcal{X}$ creates all the limits that $\mathcal{X}$ has, in particular the pullbacks, which are used to define the change-of-base functor $f^*$.
        \[
        % https://tikzcd.yichuanshen.de/#N4Igdg9gJgpgziAXAbVABwnAlgFyxMJZABgBpiBdUkANwEMAbAVxiRAB12BbOnACwBOXYAAUcAXwD6wTj34BjRsAAa4gHoAVcQAoAmqU4APLAHIAlCHGl0mXPkIoATOSq1GLNrN6DhYqTO5vRQYVdS1tZQN2YwsrG2w8AiIyAEZXemZWRA5A-iFRCWkvBSVVHV1Y6xAMBPsiZzTqDI9s4p8C-zbg0J1lWNcYKABzeCJQADMBCC4kMhAcCCQUpvcskHG1ACpLKsnppeoFpABmFcy2AFVNHYmpmcQ5o8RnN3Psq40b9bukF6fT14tdZbSwUcRAA
\begin{tikzcd}
    {\mathrm{Pt}_{\mathcal{X}^T}(Y,\xi')} \arrow[rr, "f^*"] \arrow[d, "U^T"] &  & {\mathrm{Pt}_{\mathcal{X}^T}(X,\xi)} \arrow[d, "U^T"] \\
    \mathrm{Pt}_{\mathcal{X}}(Y) \arrow[rr, "f^*"]                           &  & \mathrm{Pt}_{\mathcal{X}}(X)                         
    \end{tikzcd}
    \]
    Let $\alpha$ be a morphism in $\mathrm{Pt}_{\mathcal{X}^T}(Y,\xi')$. Suppose that $f^*(\alpha)$ in iso, we want to prove that $\alpha$ is iso. Since $f^*(\alpha)$ in iso, we have that $U^T(f^*(\alpha))$ is iso. By the commutativity of the diagram, we have that $U^T(f^*(\alpha)) = f^*(U^T(\alpha))$. Thus $f^*(U^T(\alpha))$ is an iso. Since $\mathcal{X}$ is protomodular, $U^T(\alpha)$ is an iso. Finally, $U^T$ reflects isomorphisms, so $\alpha$ is an iso.
    \end{proof}
\end{subbox}

We now proceed with results that will allow us to construct coproducts in the Eilenberg–Moore category, assuming they exist in the base category. The following will be rather intricate and technical; readers may choose to skip directly to the final outcome \Cref{algebras category has colim if the base category has coprod}.

For the sake of clarity, we first present the following technical definition.

\begin{mainbox}{}
    \begin{defi}
        Let $U:\mathcal{X} \to \mathcal{A}$ be a functor, let $X$ be an object in $\mathcal{X}$ and let $\{(f_i,g_i)\}_{i\in I}$ be a family of pair of parallel arrows in $\mathcal{A}$, where $f_i,g_i: A_i \to U(X)$. A morphism $q:X \to Q$ in $\mathcal{X}$ is said to be a $U$\textbf{-coequalizer} of the pair $\{(f_i,g_i)\}_{i\in I}$ if 
        \begin{itemize}
            \item $U(q) f_i = U(q) g_i \ \forall i \in I$ in $\mathcal{A}$;
            \item If $h:X \to Y$ is a morphism in $\mathcal{X}$ such that $U(h) f_i = U(h) g_i \ \forall i \in I$ in $\mathcal{A}$, then there exists a unique morphism $t:Q \to Y$ in $\mathcal{X}$ such that $t q = h$.
        \[
% https://tikzcd.yichuanshen.de/#N4Igdg9gJgpgziAXAbVABwnAlgFyxMJZARgBoAGAXVJADcBDAGwFcYkQBVACgA0BKEAF9S6TLnyEUAJgrU6TVu24BFAcNHY8BIuVk0GLNohABBISJAZNEojOJyDi49wCaai1fHaUAFj3zDdh5zDS9JZABWf0cjEGUQyzEtcKj7fQVYlyE5GCgAc3giUAAzACcIAFskXRAcCCQyAKdOLgBHdxLyqsQZWvrEGrgACyxinCQAWkaY9mKAfSwEssqkXrrqmmHR8cRpjPY8hZAaRnoAIxhGAAUkm2NSrDyh8fUQZe6a9cQAZnTA5y4QwEJ3Olxu1m8IAeTxeFneDRoX1+TVi3BwwJAjCwYFiUHow1ySy6SD8fSQURR7FaRJWiApXwAbH9muMTtjcfihoTXvDEKTGczYkNjpjQddbpDoc9soIgA
\begin{tikzcd}
A \arrow[r, "f_i", shift left] \arrow[r, "g_i"', shift right] & U(X) \arrow[r, "U(q)"] \arrow[rd, "U(h)"'] & U(Q) \arrow[d, "U(t)", dashed] &  & X \arrow[r, "q"] \arrow[rd, "h"'] & Q \arrow[d, "t", dashed] \\
                                                              &                                            & U(Y)                           &  &                                   & Y                       
\end{tikzcd}
        \]
        \end{itemize}
    \end{defi}
\end{mainbox}

Before continuing, recall the universal property of the unit in an adjunction. Consider a functor $U: \mathcal{A} \to \mathcal{B}$ right adjoint to a functor $F$. Then for each object $A$ in $\mathcal{A}$, $B$ in $\mathcal{B}$, and $f:B \to U(A)$ in $\mathcal{B}$ there exists a unique morphism $g: F(B) \to A$ in $\mathcal{A}$ such that $f=U(g)\eta_B$. We denote $g = \bar{f}$. 
\[
% https://tikzcd.yichuanshen.de/#N4Igdg9gJgpgziAXAbVABwnAlgFyxMJZABgBpiBdUkANwEMAbAVxiRACEQBfU9TXfIRQBGclVqMWbAKoAxABTsAlN14gM2PASKjh4+s1aIQ0+QEEVPPpsFEAzGOoGpxhctXWB2lA71PJRiBm3OIwUADm8ESgAGYAThAAtkiiIDgQSABM-oYy8uGWavFJSGRpGYipzoEAOjUwOHQA+pxWIMXJiGXpWdQMdABGMAwACvxaQiBxWOEAFjggOS7tHu0JnQ7lSAAsS4HhIVxAA
\begin{tikzcd}
B \arrow[r, "\eta_B"] \arrow[rd, "f"'] & UF(B) \arrow[d, "U(g)"] &  & F(B) \arrow[d, "g"] \\
                                       & U(A)                    &  & A                  
\end{tikzcd}
\]

Also recall that a monad $(T,\mu,\eta)$ induces a functor $U^T$ and a functor $F^T$ defined as follows:

\begin{minipage}{0.5\textwidth}
    \[
    % https://tikzcd.yichuanshen.de/#N4Igdg9gJgpgziAXAbVABwnAlgFyxMJZABgBpiBdUkANwEMAbAVxiRAB12BbOnACwDGjYAA0AvgD0AKiDGl0mXPkIoyARiq1GLNpx78hDUWNmaYUAObwioAGYAnCFyRkQOCEjXU4fLLZxIALReWsysiCAAqtKy8iAOTp7U7i7evv5BIfRhbABiMWIUYkA
\begin{tikzcd}
\mathcal{X}^T \arrow[d, "U^T", shift left] \\
\mathcal{X} \arrow[u, "F^T", shift left]  
\end{tikzcd}
    \]
\end{minipage}
\begin{minipage}{0.2\textwidth}
\[
% https://tikzcd.yichuanshen.de/#N4Igdg9gJgpgziAXAbVABwnAlgFyxMJZABgBpiBdUkANwEMAbAVxiRAAoANUgHR4A8sAShABfUuky58hFAEZyVWoxZtOYpTCgBzeEVAAzAE4QAtkjIgcEJAqt0sDNqbpo41kNXrNWiEAFUAPQAVDVEgA
\begin{tikzcd}
{(X,\xi)} \arrow[r, "U^T", maps to] & X
\end{tikzcd}
\]
\[
% https://tikzcd.yichuanshen.de/#N4Igdg9gJgpgziAXAbVABwnAlgFyxMJZABgBpiBdUkANwEMAbAVxiRAA0QBfU9TXfIRQBGclVqMWbABQAVaewCUpADoqAtkwD6S7uJhQA5vCKgAZgCcI6pGRA4ISUfbpYGbdXTRwHIavWZWRBAAMQA9WT0uIA
\begin{tikzcd}
X \arrow[r, "F^T", maps to] & {(T(X),\mu_X)}
\end{tikzcd}
\]
\end{minipage}

\begin{mainbox}{}
    \begin{lemma}\label{U-coeq}
        Let $U:\mathcal{X} \to \mathcal{A}$ be a functor right adjoint to a functor $F:\mathcal{A} \to \mathcal{X}$, let $X$ be an object in $\mathcal{X}$ and let $\{(f_i,g_i)\}_{i\in I}$ be a family of parallel arrows in $\mathcal{X}$, where $f_i,g_i: A_i \to U(X)$. Consider the family of unique morphisms induced by the universal property of the unit $\{(\bar{f}_i,\bar{g}_i)\}_{i\in I}$, such that $f_i=U(\bar{f}_i)\eta_{A_i}$ and $g_i=U(\bar{g}_i)\eta_{A_i}$. If $q:X \to Q$ is a coequalizer of the family $\{(\bar{f}_i,\bar{g}_i)\}_{i\in I}$ in $\mathcal{X}$, then $q$ is a $U^T$-coequalizer of the family $\{(f_i,g_i)\}_{i\in I}$ in $\mathcal{A}$.
       
        \begin{minipage}{0.5\textwidth}
\[
% https://tikzcd.yichuanshen.de/#N4Igdg9gJgpgziAXAbVABwnAlgFyxMJZABgBpiBdUkANwEMAbAVxiRAEEB9LEAX1PSZc+QigBM5KrUYs2AVQBiACi5YAlHwEgM2PASISxU+s1aIQcpQA0NvKTCgBzeEVAAzAE4QAtkjIgcCCQARmoTWXMAHUiYHDpOYFVeTXcvX0R-QKQJaVM2N24QajgACyw3HCQAWmD+VJ8Q6izEHPCzCyVogCM6D2A3Xm4NYrKK6tqtTwaMpqCWkfLKxFDciJBHQuoGOi6YBgAFIT1REA8sRxLKupAp9JXmnNLFxtX2y27e4EdB9SKQbd2ByOIjYZwuVwovCAA
\begin{tikzcd}
A_i \arrow[rr, "\eta_{A_i}"] \arrow[rrdd, "f_i", shift left] \arrow[rrdd, "g_i"', shift right] &  & UF(A_i) \arrow[dd, "U(\bar{f}_i)", shift left] \arrow[dd, "U(\bar{g}_i)"', shift right] \\
                                                                                               &  &                                                                                         \\
                                                                                               &  & U(X)                                                                                   
\end{tikzcd}
\]
        \end{minipage}
        \begin{minipage}{0.5\textwidth}
\[
% https://tikzcd.yichuanshen.de/#N4Igdg9gJgpgziAXAbVABwnAlgFyxMJZABgBpiBdUkANwEMAbAVxiRADEAKAQQH0sAlCAC+pdJlz5CKAIzkqtRizYANEWJAZseAkQBM86vWatEIAIoiFMKAHN4RUADMAThAC2SMiBwQkckDgACywnHCQAWgDjZTMAHTiAIzoXYCdhfnVnN09Eb19-amDQ8MRopVMQBOTU235hEGoGOkSYBgAFCR1pEBcsWyDw0WyPQp8-RANFEzYARythIA
\begin{tikzcd}
F(A_i) \arrow[r, "\bar{f}_i", shift left] \arrow[r, "\bar{g_i}"', shift right] & X \arrow[r, "q"] & Q
\end{tikzcd}
\]
\[
% https://tikzcd.yichuanshen.de/#N4Igdg9gJgpgziAXAbVABwnAlgFyxMJZABgBpiBdUkANwEMAbAVxiRAEEB9LEAX1PSZc+QigCM5KrUYs2AVQAUADQCUfASAzY8BIgCZJ1es1aIQigIpreUmFADm8IqABmAJwgBbJGRA4ISBIgcAAWWC44SAC0QcayZi7c6q4e3oi+-oHUoeGRiLEypiD2SdQMdABGMAwACkI6oiBuWPYhkfwpXll+AYgG0ibyCgCO1hS8QA
\begin{tikzcd}
A_i \arrow[r, "f_i", shift left] \arrow[r, "g_i"', shift right] & U(X) \arrow[r, "U(q)"] & U(Q)
\end{tikzcd}
\]
        \end{minipage}
    \end{lemma}
\end{mainbox}

\begin{subbox}{}
\begin{proof}
    Assume $q:X \to Q$ is a coequalizer of the family $\{(\bar{f}_i,\bar{g}_i)\}_{i\in I}$ in $\mathcal{X}$, we want to prove that $q$ is a $U^T$-coequalizer of the family $\{(f_i,g_i)\}_{i\in I}$ in $\mathcal{A}$. First we show, for each $i$, that $U(q)f_i=U(q)g_i$. Assume that $q\bar{f}_i=q\bar{g}_i$.
    \begin{align*}
        q\bar{f}_i=q\bar{g}_i &\implies U(q) U(\bar{f}_i) = U(q) U(\bar{g}_i) \\
        &\implies U(q) U(\bar{f}_i) \eta_{A_i}= U(q) U(\bar{g}_i) \eta_{A_i} \\
        &\implies U(q)f_i=U(q)g_i
    \end{align*}

    We aim to prove now the universal property of the $U$-coequalizer. Let $h:X \to Y$ be a morphism such that $U(h)f_i=U(h)g_i$. By the universal peroperty of the unit applied to $U(h)f_i=U(h)g_i : A_i \to U(Y)$ there exist a unique morphism $k:F(A_i) \to Y$ such that $U(k) \eta_{A_i}=U(h)f_i=U(h)g_i$. Both $h\bar{f}_i$ and
    $h\bar{g}_i$ satisfies such condition; in fact $U(h\bar{f}_i) \eta_{A_i}= U(h) U(\bar{f}_i) \eta_{A_i} = U(h)f_i$ and similarly for $h\bar{g}_i$. Thus by the uniqueness $h\bar{f}_i=k=h\bar{g}_i$.
    \[
% https://tikzcd.yichuanshen.de/#N4Igdg9gJgpgziAXAbVABwnAlgFyxMJZABgBpiBdUkANwEMAbAVxiRAEEB9LEAX1PSZc+QigBM5KrUYs2AVQBiACi5YAlHwEgM2PASISAjFPrNWiEHKUANDf0G6RB0mJMzzlpQE07WncP0UABZJalNZC2VVXwcA0WQQ4zD3NmtNWL14kNdkszYvPikYKABzeCJQADMAJwgAWyRDahwIJAkQOAALLEqcJABaJuk8iysAHTGAIzpq4ErebhiQGvqkMhAWttyI5e4Qai6evsRB+2XahsR1zcR2w97G7Y8SveoGOkmYBgAFIUy2apYEqdPpnFaXa6tRAAZmonzAUAG0PW92O63CHisnTUlW4AF4sWoXjw3h8vr9HIEQIDgaCtOCthsobDhjtCelzqsrs0oUMMWwJjAcHROMBVLwOQzEEMbndug9paTPj8-k4LDSQftWZilBNprMSgt1JKLo8mUgWfDEYgAGwAdie8iUAGsllKQubEABWA7y46DR0WPUzYCGvZg03enlIG2BkCdE1cj03H0dP1m-lBqYh+avEDvZUUuIAoGaiNJ6O2uEwBEx9EpCzOwq8IA
\begin{tikzcd}
A_i \arrow[rrd, "f_i", shift left] \arrow[rrd, "g_i"', shift right] \arrow[rrdd, "U(h)f_i=U(h)g_i"', bend right] \arrow[rr, "\eta_{A_i}"] &  & UF(A_i) \arrow[d, "U(\bar{f}_i)", shift left] \arrow[d, "U(\bar{g}_i)"', shift right] \arrow[dd, "U(k)", bend left=67] &  & F(A_i) \arrow[d, "\bar{g}_i", shift left] \arrow[d, "\bar{f}_i"', shift right] \arrow[dd, "k", bend left=60] \\
                                                                                                                                          &  & U(X) \arrow[d, "U(h)"]                                                                                                 &  & X \arrow[d, "h"]                                                                                             \\
                                                                                                                                          &  & U(Y)                                                                                                                   &  & Y                                                                                                           
\end{tikzcd}
    \]
    Then, by the universal property of the coequalizer $q$ in $\mathcal{X}^T$, there exists a unique morphism $t:Q \to Y$ sucht that $tq=h$.
\end{proof}
\end{subbox}

\begin{mainbox}{}
\begin{prop}\label{coeq for the three diagrams}
    Let $\mathcal{X}$ be a category, let $(T,\mu,\eta)$ be a monad on $\mathcal{X}$ and let $\{(A_i,\xi_i)\}_{i\in I}$ be a family of objects in $\mathcal{X}^T$. Assume that the coproducts $\amalg_iA_i$ and $\amalg_iT(A_i)$ exist in $\mathcal{X}$ with respective canonical injections $\{\iota_i\}_{i\in I}$ and $\{\iota'_i\}_{i\in I}$. If the following parallel arrows in 
    $\mathcal{X}^T$ admits a coequalizer $q:T(\amalg_i)\to (Q,\xi_Q)$
    \begin{equation}\label{diagr 1 coproducts}
% https://tikzcd.yichuanshen.de/#N4Igdg9gJgpgziAXAbVABwnAlgFyxMJZABgBoBGAXVJADcBDAGwFcYkQAdDgISwHMAFABUuAI34CucZgFsA+liECAggoCUY-mtJcZzLrz5qQAX1LpMufIRTlSxanSat2BiUIB6AJkkdp8rFUsbV19Hi1TcxAMbDwCIi8KRwYWNkROcMElKVkFIJCOPTcjU0cYKD54IlAAMwAnCBkkMhAcCCREkFEYMChmmhSXdOy-XKwuAA8sdRAaRnpuxgAFSzibEDr+AAscSNqGpsQWtqQ7J1T2IWQuRigIHDhSEfwcenUdDlv7uEpZrp6+kczPtGqcaCdEJ1BmkMno5MAcgEgiY-t1es0TJQTEA
\begin{tikzcd}
                                                                                                                  & {\Big(T^2(\amalg_iA_i),\mu\Big)} \arrow[rd, "\mu_{\amalg_iA_i}"] &                              \\
{\Big(T\big(\amalg_iT(A_i)\big),\mu\Big)} \arrow[rr, "T(\amalg_i\xi_i)"'] \arrow[ru, "{T[\ldots,T(\iota_i),\ldots]}"] &                                                              & {\Big(T(\amalg_iA_i),\mu\Big)}
\end{tikzcd}
\end{equation}
    then both the following diagrams of parallel arrow in $\mathcal{X}$ have $p$ as $U^T$-coequalizer $\forall j \in I$.
\begin{equation}\label{diagr 2 coproducts}
    % https://tikzcd.yichuanshen.de/#N4Igdg9gJgpgziAXAbVABwnAlgFyxMJZABgBpiBdUkANwEMAbAVxiRAB124mBbAfSwAVABQBBAQEoQAX1LpMufIRQAmclVqMWbEZ278s4rFNnzseAkQCMpKxvrNWiDl14CjMjTCgBzeEVAAMwAnCB4kMhAcCCQbTUc2ZE4GKAgcOFJddnwcOklSZNT0ihk5EBCwiOpopDUQACMYMCgq+O1nPTcsAAJOAA8sARBqBjpGhgAFBQtlEGCsHwALHFKg0PDEOprEOMbm1od2lxhcvmBOgyNpYZBR8anzJTZ5pZXpCmkgA
\begin{tikzcd}
\amalg_iT(A_i) \arrow[rr, "{[\ldots,T(\iota_i),\ldots]}"] \arrow[rd, "\amalg_i \xi_i"'] &                                           & T(\amalg_iA_i) \\
                                                                                    & \amalg_iA_i \arrow[ru, "\eta_{\amalg_iA_i}"'] &             
\end{tikzcd}
\end{equation}

\begin{equation}\label{diagr 3 coproducts}
% https://tikzcd.yichuanshen.de/#N4Igdg9gJgpgziAXAbVABwnAlgFyxMJZABgBpiBdUkANwEMAbAVxiRABUAKAQQH0ArAJQgAvqXSZc+QigDM5KrUYs2XADpq6AW0YBzXlgAEfLMLETseAkQCMpG4vrNWiEH36jxIDJelEATPaOyi4gGtp6BiaiijBQuvBEoABmAE4QWkhkIDgQSHZKzqqcGvg4dAJmXmkZWdS5SIGFKq4aAB5YvMD8IiDUAEYwYFB1IAx0gwwACpJWMiCpWLoAFjieKemZiE0NiPLNoaUQ5QLrIDVb+7sFTi1hajAnwOE6DPpGJr0DQyOI2eOTGa+ayuRYrNYiCgiIA
\begin{tikzcd}
T(A_j) \arrow[rrr, "T(\iota_j)"] \arrow[rd, "\xi_{j}"'] &                          &                                                & T(\amalg_i A_i) \\
                                                        & A_j \arrow[r, "\iota_j"] & \amalg_iA_i \arrow[ru, "\eta_{\amalg_i A_i}"'] &                
\end{tikzcd}
\end{equation}
\end{prop}
\end{mainbox}

\begin{subbox}{}
    \begin{proof}
        Consider the diagram \eqref{diagr 1 coproducts} of parallel arrows in $\mathcal{X}^T$, and $p$ its coequalizer, here below explicitly.
        \[
    % https://tikzcd.yichuanshen.de/#N4Igdg9gJgpgziAXAbVABwnAlgFyxMJZABgBoBmAXVJADcBDAGwFcYkQAVAHS4CMsA5gAoecZgFsA+lg5CAgtICUPfgMUgAvqXSZc+QigBMpQ9TpNW7DgD1DIrmKlYFWdVp3Y8BIgBYKZhhY2RE57R2kXN20QDE99ImNiAItgzmtyMIkIpU1o2L1vFDIARmSgq1sVQUynWUiqtVyPAoNkP1KaQMsQmztRLOccjTMYKAF4IlAAMwAnCHEkMhAcCCRjEF4YMChFzpSrGukeAA8sHJpGek3GAAVdLwMQGcEACxwmkFn53eXVxGK9uUesgeIwoBAcHBSLIePgcPQlKRQeDIZQPl8Fv8aCs1oDuiAeOJmJJgP0nC4NOi5pjyNi-gDzECCVwiSSYQ4BpFKe5PtSkH5fj9Ltc7nFCk9Xu88alCcTSRzavIlNzohikABWOm4xn42UksnZLAq6Z8xACnGIWk61K9EFcMEQqHsuEI1xI+0ouBonlqy1axCa60HPXy8KDI1RE3fM3+wNdG2HLAnM6uTSUDRAA
\begin{tikzcd}
                                                                                                                                          &  & T^3(\amalg_iA_i) \arrow[dd, "\mu_{T(\amalg_iA_i)}", pos=0.8] \arrow[rrd, "T\big(\mu_{\amalg_iA_i}\big)"] &  &                                              \\
T^2\big(\amalg_iT(A_i)\big) \arrow[dd, "\mu_{\amalg_iT(A_i)}"'] \arrow[rru, "{T^2[\ldots,T(\iota_i),\ldots]}"] \arrow[rrrr, "T(\amalg_i\xi_i)", pos=0.3, crossing over] &  &                                                                                   &  & T^2(\amalg_iA_i) \arrow[dd, "\mu_{\amalg_iA_i}"] \\
                                                                                                                                          &  & T^2(\amalg_iA_i) \arrow[rrd, "\mu_{\amalg_iA_i}"]                                     &  &                                              \\
T\big(\amalg_iT(A_i)\big) \arrow[rrrr, "T(\amalg_i\xi_i)"'] \arrow[rru, "{T[\ldots,T(\iota_i),\ldots]}"]                                      &  &                                                                                   &  & T(\amalg_iA_i)                                
\end{tikzcd}
\]
The top and bottom arrows from $(T(\amalg_iT(A_i)),\mu)$ to $(T(\amalg_iA_i),\mu)$ of the diagram \eqref{diagr 1 coproducts} correspond, by the universal property of the unit of the adjunction, respectively to the top and bottom arrows of the diagram \eqref{diagr 2 coproducts}. In order to prove so, by the uniqueness of the morphism induced by the universal property of the unit, it suffies to show that the image of the diagram \eqref{diagr 1 coproducts} through $U^T$, composed with the unit, is the diagram \eqref{diagr 2 coproducts}. Such composition results to the following diagram.
\[
% https://tikzcd.yichuanshen.de/#N4Igdg9gJgpgziAXAbVABwnAlgFyxMJZAJgBpiBdUkANwEMAbAVxiRABUAdTgIywHMAFNzhMAtgH0s7QQEEpASm59+CkAF9S6TLnyEUAZlIAGKrUYs27AHrFhnUZKzysazdux4CRACykAjGb0zKyIHPaOUi5uWiAYnnpE-iZBFqEcygIR4lIy0ZmqGrHxut4oxgGpIWwiOdJyikUepfokldTBlmERaHSworhuZjBQ-PBEoABmAE4QYkgVIDgQSGQgPDBgUAsdaVbZTtwAHliN1Ax0GwwACjpe+iDTAgAWOE0gM3NIyUsriGuddLcMRMCTAWpOFzqd6feaIIy-b67aphdjIbgMKAQHBwUgybj4HB0RSkDFYnEUGGzOF+RHw5FdEDcGBEsEQ3INACSCmh7g+1KQtOWO3MKKZnBZxPBDjqeQk3Oh50uMBud0SYSe-FeGgo6iAA
\begin{tikzcd}
                                                                                   & T\big(\amalg_iT(A_i)\big) \arrow[rr, "{T[\ldots,T(\iota_i),\ldots]}"] &                                                        & T^2(\amalg_iA_i) \arrow[rd, "\mu_{\amalg_iA_i}"] &              \\
\amalg_iT(A_i) \arrow[ru, "\eta_{\amalg_iT(A_I)}"] \arrow[rrd, "\eta_{\amalg_iT(A_I)}"'] &                                                                     & (\spadesuit)                                           &                                              & T(\amalg_iA_i) \\
                                                                                   &                                                                     & T\big(\amalg_iT(A_i)\big) \arrow[rru, "T(\amalg_i\xi_i)"'] &                                              &             
\end{tikzcd}
\]
Consider now the top arrow in the diagram $(\spadesuit)$ from $\amalg_iT(A_i)$ to $T(\amalg_iA_i)$. Using the naturality of $\eta$, we have that
\[
% https://tikzcd.yichuanshen.de/#N4Igdg9gJgpgziAXAbVABwnAlgFyxMJZABgBpiBdUkANwEMAbAVxiRAB124mBbAfSwAVABQBBAQEoQAX1LpMufIRRkAjFVqMWbEZ278s4rFNnzseAkQBM5DfWatEIQZwBGWAObC9vASKMSbp4mciAY5krWpOrU9tpOggB6Vt5cvoaSMhowUB7wRKAAZgBOEDxIZCA4EEiqsVqOIMicDFAQOHCkuuz4OHSSpC1tHRQg1Ax0rjAMAAoKFsogxZ4AFjgyoSVlFdTVSDaaDmycMH18wD4G-pLSG0Wl5Yh1VTWIAMz1R04nZ8Dd+gIArdTCAto8Dnt3p94s5muxWu1Ot1ev1jIN4cM4KNpBRpEA
\begin{tikzcd}
\amalg_iT(A_i) \arrow[d, "{[\ldots,T(\iota_i),\ldots]}"'] \arrow[rr, "\eta_{\amalg_iT(A_i)}"] &  & T\big(\amalg_iT(A_i)\big) \arrow[d, "{T[\ldots,T(\iota_i),\ldots]}"] \\
T(\amalg_iA_i) \arrow[rr, "\eta_{T(\amalg_iA_i)}"]                                            &  & T^2(\amalg_iA_i)                                                    
\end{tikzcd}
\]
\begin{align*}
    \mu_{\amalg_iA_i} T[\ldots,T(\iota_i),\ldots] \eta_{\amalg_iT(A_i)} &= \mu_{\amalg_iA_i} \eta_{T(\amalg_iA_i)} [\ldots,T(\iota_i),\ldots] \\[1.5ex]
    &= 1_{T(\amalg_iA_i)} [\ldots,T(\iota_i),\ldots] \\[1.5ex]
    &= [\ldots,T(\iota_i),\ldots] 
\end{align*}
which is the top arrow in \eqref{diagr 2 coproducts}. Then, for the bottom arrow from $\amalg_iT(A_i)$ to $T(\amalg_iA_i)$ in $(\spadesuit)$, by the naturality of $\eta$ and the property of the monad, we have that
\[
% https://tikzcd.yichuanshen.de/#N4Igdg9gJgpgziAXAbVABwnAlgFyxMJZABgBpiBdUkANwEMAbAVxiRAB124mBbAfSwAVABQBBAQEoQAX1LpMufIRQAmclVqMWbQZwBGWAObDO3fkLGT9RqbPnY8BImoCMG+s1aIQI07wHiWLZyIBgOSkRkbtQe2t5+5oEyGjBQhvBEoABmAE4QPEhkIDgQSC4xWl4c7DA4dHzACQIigRLSMiG5+YXUJUgAzBWebE1YnAAeWAIg1Ax0ejAMAAoKjsogOUYAFjgd2XkFiIPFpYhqmsPxNXUNo4HtdiBdh+UnSOexVb5c-mPsk5JktIgA
\begin{tikzcd}
\amalg_iT(A_i) \arrow[rr, "\eta_{\amalg_iT(A_i)}"] \arrow[d, "\amalg_i\xi_i"'] &  & T\big(\amalg_iT(A_i)\big) \arrow[d, "T(\amalg_i\xi_i)"] \\
\amalg_iA_i \arrow[rr, "\eta_{\amalg_iA_i}"]                                 &  & T(\amalg_iA_i)                                       
\end{tikzcd}
\]
\[
T(\amalg_i\xi_i)\eta_{\amalg_iT(A_i)} = \eta_{\amalg_iA_i}(\amalg_i\xi_i)
\]
which is the bottom arrow in \eqref{diagr 2 coproducts}. By \Cref{U-coeq}, it follows that the diagram \eqref{diagr 2 coproducts} has $p$ as $U^T$-coequalizer.

It remains to prove that the diagram \eqref{diagr 3 coproducts} has $p$ as $Us^T$-coequalizer. We already know that \eqref{diagr 2 coproducts} has $p$ as $U^T$-coequalizer. Note that these two diagrams are related by the canonical injection $\iota'_i:T(A_i) \to \amalg_iT(A_i)$, as shown in the diagram. In particular, by the universal property of the coproduct $(\amalg_i \xi_i) \iota'_j = \iota_j\xi_j$ and $[\ldots,T(\iota_i),\ldots] \iota'_j = T(\iota_j)$.
\[
% https://tikzcd.yichuanshen.de/#N4Igdg9gJgpgziAXAbVABwnAlgFyxMJZARgBoAGAXVJADcBDAGwFcYkQAdDuZgWwH0sAFQAUAQUEBKEAF9S6TLnyEUAZgrU6TVu1FceArBKzS5C7HgJEATKWKaGLNok7c+g47PkgMF5UXINGkcdF1EJACtTTRgoAHN4IlAAMwAnCF4kQJAcCCQyLSd2ZC5GKAgcOFI9DnwceilSUvLKyi8U9MzEbNykWxAAIxgwKCzg7WdXA0EAAi4ADyxBEBpGeiHGAAVFSxUQVKw4gAscdpA0jL6aXsQCoZGxwtDXGHr+YH13I0EZFZA1jbbPxWFwHY6nMznTpIdQ5PLdcZFFxcOr0ADk-AiZwuXVhNzuw1GiAALABORHPGqozGmbw4mHXeH9e5EgC0qmyIUmKIqDQiCyWWNW6xgWx2-lBhxOskoMiAA
\begin{tikzcd}
T(A_j) \arrow[r, "\iota'_j"] \arrow[rrr, "T(\iota_j)", bend left=35] \arrow[rrd, "\iota_j\xi_j"', bend right] & \amalg_iT(A_i) \arrow[rr, "{[\ldots,T(\iota_i),\ldots]}"] \arrow[rd, "\amalg_i \xi_i"'] &                                           & T(\amalg_iA_i) \\
                                                                                                              &                                                                                     & \amalg_iA_i \arrow[ru, "\eta_{\amalg_iA_i}"'] &             
\end{tikzcd}
\]
The first condition of the $U^T$-coequalizer for \eqref{diagr 3 coproducts} is obviously satisfied, just composing with $\iota'_j$ the condition of the $U^T$-coequalizer of \eqref{diagr 2 coproducts}. 
\begin{align*}
    U^T(q) &[\ldots,T(\iota_i),\ldots] = U^T(q) \eta_{\amalg_iA_i} \amalg_i \xi_i \implies \\[1.5ex]
    &\implies U^T(q) [\ldots,T(\iota_i),\ldots] \iota'_j= U^T(q) \eta_{\amalg_iA_i} \amalg_i \xi_i \iota'_j \\[1.5ex]
    &\implies U^T(q) T(\iota_j) = U^T(q) \iota_j \xi_j 
\end{align*}

The universal property of the $U^T$-coequalizer of \eqref{diagr 3 coproducts} also holds. In fact, let $h$ is a morphism in $\mathcal{X}^T$ such that $U^T(h) T(\iota_j) = U^T(h) \iota_j \xi_j \ \ \forall j \in I$, then 
\begin{align*}
    U^T(h) &T(\iota_j) = U^T(h) \iota_j \xi_j \ \ \ \forall j \in I \implies \\[1.5ex]
    &\implies U^T(h) [\ldots,T(\iota_i),\ldots] \iota'_j= U^T(h) \eta_{\amalg_iA_i} \amalg_i \xi_i \iota'_j \ \ \ \forall j \in I \\[1.5ex]
    &\implies U^T(h) [\ldots,T(\iota_i),\ldots] = U^T(h) \eta_{\amalg_iA_i} \amalg_i \xi_i
\end{align*}
and by the universal property of the $U^T$-coequalizer of \eqref{diagr 2 coproducts} there exists a unique morphism $k$ such that $kq=h$.
\end{proof}
\end{subbox}

We are now ready to prove the result that will allow us to construct coproducts in the Eilenberg–Moore category, assuming they exist in the base category. This is possible if we assume the existence of reflexive coequalizers in the Eilenberg–Moore category.

\begin{mainbox}{}
    \begin{prop}\label{algebras category has colim if the base category has coprod}
        Let $T:\mathcal{\mathcal{X}} \to \mathcal{X}$ be a monad. If $\mathcal{X}$ has finite coproducts and the Eilenberg-Moore category $\mathcal{X}^T$ has reflexive coequalizers, then $\mathcal{X}^T$ has finite coproducts. Moreover $\mathcal{X}^T$ is finitely cocomplete.
    \end{prop}
\end{mainbox}

\begin{subbox}{}
    \begin{proof}
        Let $(T,\mu,\eta)$ be the given monad, and suppose that $\mathcal{X}$ has finite coproducts. Let $(A,\xi_A)$ and $(B,\xi_B)$ be two objects in $\mathcal{X}^T$. We aim to prove that $(A,\xi_A)$ and $(B,\xi_B)$ have a coproduct in $\mathcal{X}^T$. All the coproducts in the diagram below are taken in $\mathcal{X}$, and therefore they exist. Let $\iota_1$ and $\iota_2$ be the canonical injections in $\mathcal{X}$ of $A$ and $B$ into $A+B$, respectively. Our goal is to prove that the diagram below is composed by a reflexive pair in $\mathcal{X}^T$. Then, by assumption, it admits a coequalizer in $\mathcal{X}^T$, denoted by $q$. This diagram correspond to the diagram \eqref{diagr 1 coproducts} in \Cref{coeq for the three diagrams}.
            \[
% https://tikzcd.yichuanshen.de/#N4Igdg9gJgpgziAXAbVABwnAlgFyxMJZABgBoAWAXVJADcBDAGwFcYkQAVAHS4CMsA5owAUAAg7CAggEoA1BIBC0nvwEAnaSAC+pdJlz5CKAEwVqdJq3YTJspdt0gM2PASLkzNBizaIQARQc9F0MiMgBGc28rPw4APWMVQRFxKTlFZT5BDSCnfVcjZFNIr0tfTgSpO00dYIM3FA8Six9rYX8ax2d6wvDSAGYosutK23tavJCG5D7iIdbYuP6q8fMYKAF4IlAAMzUIAFskMhAcCCQ+kDgACywdnAvShc5kHgAPLAB9SVJ3r4VKCAaIx6LwYIwAAr5UJ+NSCa4PCZ7Q6PU7nRAnaLlCQ8GA4ejfWS4-GfezA0HgqFTIwgLBgbCwIFXW73C5I-ZHRCXM5IUwtGIgACOuWRnP6NB5GKeAp4B2Yn2ANnSwiUWiZILBkOhDRAcIECJFHKQHjRqKx7Fl8uAYzV7JRiAArBL0XzzX4-p9AnbOSbJeL+djKsSCbZg6TNOTNVSeuw9QaaDc7g8ud6kP7JSbE6yudLA8ZXlwPt9foX-oDI5TtTS44jHKLjc6kE6A21BZ1dkapabEAA2Cta6nsRgwVm56zIHFcfAk8LSUiT6cE4zSQGp3uNnMt91cOUKm2G+0Adg3fa3IEtCps1Vtdc76fRx7P8XzC4gM7nr5Jy9Xt6PG5NbqcMIF7WnYWg1JQWhAA
\begin{tikzcd}
                                                                                                                                                    & T^3(A+B) \arrow[ddd, "\mu_{T(A+B)}", pos=0.75] \arrow[rd, "T(\mu_{A+B})"] &                                                                                                     &  &                           \\
T^2\bigl( T(A)+T(B)\bigr) \arrow[ddd, "\mu_{T(A)+T(B)}"'] \arrow[rr, "{T^2(\xi_A+\xi_B)}"', shift right, crossing over, pos=0.25] \arrow[ru, "{T^2[T(\iota_1),T(\iota_2)]}"] &                                                                 & T^2(A+B) \arrow[ddd, "\mu_{A+B}"] \arrow[ll, "T^2(\eta_A+\eta_B)"', shift right, crossing over, pos=0.75] \arrow[rr, "T(q)"] &  & T(Q) \arrow[ddd, "\xi_Q"] \\
                                                                                                                                                    &                                                                 &                                                                                                     &  &                           \\
                                                                                                                                                    & T^2(A+B) \arrow[rd, "\mu_{A+B}"]                                &                                                                                                     &  &                           \\
T\bigl( T(A)+T(B)\bigr) \arrow[rr, "{T(\xi_A+\xi_B)}"', shift right] \arrow[ru, "{T[T(\iota_1),T(\iota_2)]}", pos=0.65]                                       &                                                                 & T(A+B) \arrow[ll, "T(\eta_A+\eta_B)"', shift right] \arrow[rr, "q"]                      &  & Q                        
\end{tikzcd}
\]
        The parallel arrows in $\mathcal{X}^T$ to be coequalized are
        \begin{itemize}
            \item $T[\xi_A,\xi_B]$
            \item $\mu_{A+B}T[T(\iota_1),T(\iota_2)]$
        \end{itemize}
        We proceed to show that they form a reflexive pair by exhibiting a common section. Consider $T(\eta_A+\eta_B)$; it is a common section, as shown below. It follows the first composition
        \[T[\xi_A,\xi_B]T(\eta_A+\eta_B)=T\bigl([\xi_A,\xi_B](\eta_A+\eta_B)\bigr)=T(1_{A+B})=1_{T(A+B)} \]
        and the second composition, using the naturality of $\eta$.
        \begin{align*}
        \mu_{A+B}T[T(\iota_1),T(\iota_2)]T&(\eta_A+\eta_B)=\mu_{A+B}T\bigl([T(\iota_1),T(\iota_2)](\eta_A+\eta_B)\bigr) \\[1.5ex]
        &= \mu_{A+B}T[T(\iota_1)\eta_A,T(\iota_2)\eta_B] \\[1.5ex]
        &= \mu_{A+B}T[\eta_{A+B}\iota_1,\eta_{A+B}\iota_2] \\[1.5ex]
        &= \mu_{A+B}T(\eta_{A+B} [\iota_1,\iota_2]) \\[1.5ex]
        &= \mu_{A+B}T(\eta_{A+B} 1_{A+B}) \\[1.5ex]
        &= \mu_{A+B}T(\eta_{A+B}) = 1_{T(A+B)}
        \end{align*}
        Thus it is a reflexive pair. 
        
        Let $q:T(A+B) \to Q$ be the coequalizer in $\mathcal{X}$, and $\xi_Q$ be the algebra structure of $Q$. We aim to prove that $(Q,\xi_Q)$ is the coproduct of $(A,\xi_A)$ and $(B,\xi_B)$ in $\mathcal{X}^T$. First we need to exhibit the canonical injections $\iota_{Q,1}$ and $\iota_{Q,2}$. They are respectively
        
        \begin{minipage}{\textwidth}
        \[
% https://tikzcd.yichuanshen.de/#N4Igdg9gJgpgziAXAbVABwnAlgFyxMJZABgBpiBdUkANwEMAbAVxiRAEEQBfU9TXfIRQAmclVqMWbdgGoAQt14gM2PASIAWMdXrNWiEABUAFLLkBKRX1WCiANm0S9bAIrdxMKAHN4RUADMAJwgAWyQyEBwIJABGHUl9EAAdJPwcOgB9GKsQINDY6iikUScpAxSYdIzgMy4cvLDEEqLEAGZ45wMAR3cuIA
\begin{tikzcd}
A \arrow[rr, "\iota_1"] &  & A+B \arrow[rr, "\eta_{A+B}"] &  & T(A+B) \arrow[rr, "q"] &  & Q
\end{tikzcd}
\]
\[
% As the one above 
\begin{tikzcd}
A \arrow[rr, "\iota_2"] &  & A+B \arrow[rr, "\eta_{A+B}"] &  & T(A+B) \arrow[rr, "q"] &  & Q
\end{tikzcd}
\]
\end{minipage}
        We still need to prove that they are morphisms in $\mathcal{X}^T$. We prove the first one, the second one is similar. Note that $\iota_1$ and $\eta_{A+B}$ are not morphisms in $\mathcal{X}^T$. This is not a problem, however, since we will show that the whole composition $q\eta_{A+B}\iota_1$ is a morphism in $\mathcal{X}^T$. Consider the following diagram. All diagrams commute except the left and the external ones. We aim to prove that the left diagram commutes, so that the external one commutes as well. This holds by \Cref{coeq for the three diagrams}, since it corresponds to the diagram \eqref{diagr 3 coproducts}.
        \[
        % https://tikzcd.yichuanshen.de/#N4Igdg9gJgpgziAXAbVABwnAlgFyxMJZABgBoAmAXVJADcBDAGwFcYkQBBEAX1PU1z5CKchWp0mrdhwDUAIR58QGbHgJEALGJoMWbRCAAqACllyAlIv6qhRAGzaJe9gEUrygWuElSxcbqkDEw5LXmtBdRFff0l9I1N5UKUVCO8tPx1Y9kMAPXIEi3cUr3tozOcg4xck8JKULQBGGIr4s1DxGCgAc3giUAAzACcIAFskMhAcCCQG8sCQAB0F-Bx6AH0G9yHRmZoppFEneaWYVbXgM24t4bHEQ-3EAGY5uIBHa53ELUnpxAmAuJLAAeWDWXDCIG2t2+DwArC9ssYlit1g0apCbkh4T8kA4jnETCczhd5Nx0VDcXtfgAOBEGJYjZjnS4fW7Yh60kCMegAIxgjAACp5bAZBlgugALHAgOkgBrnYKJK4QimITkPZ749jvFWYxB4h4AdllJle5L1xpxT1lwNBbm4lG4QA
\begin{tikzcd}
T(A) \arrow[dd, "\xi_A"] \arrow[rr, "T(\iota_1)"] &  & T(A+B) \arrow[rr, "T(\eta_{A+B})"] \arrow[rrd, equal] &  & T^2(A+B) \arrow[d, "\mu_{A+B}"] \arrow[rr, "T(q)"] &  & T(Q) \arrow[dd, "\xi_Q"] \\
                                                  &  &                                                               &  & T(A+B) \arrow[rrd, "q"]                            &  &                          \\
A \arrow[rr, "\iota_1"]                           &  & A+B \arrow[rr, "\eta_{A+B}"]                                  &  & T(A+B) \arrow[rr, "q"]                             &  & Q                       
\end{tikzcd}
        \]

        Once we have the canonical injections, we need to show that they satisfy the universal property of coproducts in $\mathcal{X}^T$. Let $f_i:(A_i,\xi_i) \to (K,\xi_K)$ for $i=1,2$ in $\mathcal{X}^T$. We aim to prove that there exists a unique morphism $h:(Q,\xi_Q) \to (K,\xi_K)$ in $\mathcal{X}^T$ such that $h\iota_{Q,i} = f_i$ for $i=1,2$. First, on can show that there exists a unique morphism $g=\xi_K T[\ldots,f_i,\ldots]$ in $\mathcal{X}^T$ such that $g\eta_{A+B}\iota_i=f_i$ for $i=1,2$.
        \[
% https://tikzcd.yichuanshen.de/#N4Igdg9gJgpgziAXAbVABwnAlgFyxMJZAJgBoBmAXVJADcBDAGwFcYkQAVACgEEBqAEIBKEAF9S6TLnyEU5UgBZqdJq3bcA0iPGTseAkQWkArMoYs2iEBrESQGPTKJkAjGdWXOAPWK9B2uwdpAzlSYncLdR8uLVtdYNlkIyoaczUrTQD4-USABgoI9JAeOPspHKJ8t1SPdV4ssscQpLDCzwBFUqCKlCNctrr27WUYKABzeCJQADMAJwgAWyR8kBwIJBcayIzkAB1dxigIHDhSaYB9LFJ9w+O4ShAaRnoAIxhGAAVypytZrDGABY4UpzRYbGhrJBkFTbED7AAeWHONier3eXyashAf0BwJ0IFBS0QRlW60Qxi2RX2C2YyJB8yJJMhiE2MKpuxpdPxhKQFNJUMpngRSJs3IZSHk-OJgqixD2ByOJzOl2uCruDzFYMQADYIWSVs83p9viFsf8gY82ULdjAcPRzsB+AJRPt8Hbzi56VqAOx6iUyjJcfa2+2OwQu3Zu+0uBo8xC+qW6q3sYXnEqoo0YhLsHEWzVEpPM6FvMBQZYBgkey2G9EmrG5vF2OMJ5l8mvGzE583AmglsuIFZpTzcC4xr1ElbMgAcFYAjuP-VKAJwV7iz2PixArqUz5NWVOdfNIbet3swUtIBTbod1AEbrW7otni-E6+1KwAsSUURAA
\begin{tikzcd}
                                                                                       &  &                                                                     &                                                & T(Q) \arrow[dd, "\xi_Q"] \\
T(A) \arrow[rr, "T(\eta_{A+B}\iota_1)"] \arrow[dd, "\xi_A"'] &  & T^2(A+B) \arrow[rd, "{T^2[\ldots,f_i,\ldots]}"] \arrow[rru, "T(q)"] &                                                &                                                            \\
                                                                                       &  &                                                                     & T^2(K) & Q \arrow[ddd, "h", bend left=49]                           \\
A \arrow[rr, "\eta_{A+B}\iota_1"'] \arrow[rrrrdd, "f_1"']                              &  & T(A+B) \arrow[rd, "{T[\ldots,f_i,\ldots]}", pos=0.3] \arrow[rru, "q"', pos=0.2]       &                                                & T(K) \arrow[from=uuu, "T(h)", bend left=49, crossing over] \arrow[from=lu, "\mu_K", pos=0.75, crossing over]  \arrow[from=lllluu, "T(f_1)"', crossing over] \arrow[dd, "\xi_K"]                                   \\
                                                                                       &  &                                                                     & T(K)   \arrow[rd, "\xi_K"]                      &                                                            \\
                                                                                       &  &                                                                     &                                                & K                                                         
\end{tikzcd}
        \]
        \[
        % https://tikzcd.yichuanshen.de/#N4Igdg9gJgpgziAXAbVABwnAlgFyxMJZAZgBpiBdUkANwEMAbAVxiRABUAKAQQGoAhAJQgAvqXSZc+QigAspWVVqMWbLgGlhYidjwEiAVlIGl9Zq0Qh1o8SAy7pRMgEZTKixwB6AJh4CttvZS+nKk3m7maj6cmjY6wTLIRpTUZqqWGgHxeonO5BHpINxxdpI5RHmuqe5qPFmlDiFJYQUeAIolQeUoRgAMrbVt9V2OKL0t1ZEZnEKdZaPI44qThfyiSjBQAObwRKAAZgBOEAC2SOMgOBBIecpTHMj7APp5z+Eg1Ax0AEYwDAAK8xCIEOWC2AAscCUjqcbtQrkhvCsPAAdFEADywT2snx+f0BjRkILBkOhxzOiHkl2uiCMd0KaJOTGxZNhlPhNNuaVRKKZLO0IBhFLpCMQSPpPMx-NsQqQZGpSCp3Ki3mQaIYUAgODgpGeWFI6s12oorIpADYOedcb8AUCiaCIVDkWw0TAcHQnsA+PwRGj8O6XqakAB2S2IeXK6augNegS+lH+j3OeqyxChhWIC0Sl0YrHFa34u1sB2kgWprOi8W-MBQK3ZyzPZwfEBfG0EhLFklQsvkkNhumtwuEzuO5vV2uIC6RjicRsp3uTsMADmdlgAjkHw2GAJyrmdr+ds3cZlf1kBoqUdHtH-vUceK4-Trjgw8U0+Vu8wGsPvfgzfOC5RTpJ9ZyebxXxuW4PzPN5m0HW1h0sEtuxlBcALDZwuRqSwLyxNZrwpdCMwjbCZ2jD1Yx9P0tQ9cC4LxBCO0sBgYH2bsKBEIA
\begin{tikzcd}
                                                                                       &                                                                                        &  &                                                                     &                                                & T(Q) \arrow[dd, "\xi_Q"] \arrow[ddd, "T(h)", bend left=49] \\
                                                                                       & T(A) \arrow[rr, "T(\eta_{A+B}\iota_1)"] \arrow[dd, "\xi_A"'] \arrow[rrrrdd, "T(f_1)"'] &  & T^2(A+B) \arrow[rd, "{T^2[\ldots,f_i,\ldots]}"] \arrow[rru, "T(q)"] &                                                &                                                            \\
T(B) \arrow[rrrrrd, "T(f_2)"] \arrow[dd, "\xi_B"] \arrow[rrru, "T(\eta_{A+B}\iota_2)"] &                                                                                        &  &                                                                     & T^2(K) \arrow[rd, "\mu_K"] \arrow[dd, "\mu_K"] & Q \arrow[ddd, "h", bend left=49]                           \\
                                                                                       & A \arrow[rr, "\eta_{A+B}\iota_1"'] \arrow[rrrrdd, "f_1"']                              &  & T(A+B) \arrow[rd, "{T[f_1,f_2]}"'] \arrow[rru, "q"]                 &                                                & T(K) \arrow[dd, "\xi_K"]                                   \\
B \arrow[rrrrrd, "f_2"']                                                               &                                                                                        &  &                                                                     & T(K) \arrow[rd, "\xi_K"']                      &                                                            \\
                                                                                       &                                                                                        &  &                                                                     &                                                & K                                                         
\end{tikzcd}
        \]
        Then, since $q$ $U^T$-coequalizes the same pair of arrow from $T(A)$ to $T(A+B)$ as in \eqref{diagr 3 coproducts} that have equal composition with $U^T(g)$, it follows that there exists a unique morphism $h$ in $\mathcal{X}^T$ such that $h q = g$. Thus, composing with $\eta_{A+B}\iota_i$, there exists a unique morphism $h$ in $\mathcal{X}^T$ such that $h\iota_{Q,i} = f_i$ for $i=1,2$. So the universal property of coproducts in $\mathcal{X}^T$ is satisfied.

        In order to prove that also the finite coproducts exists, it suffies to do the same proof for the initial object, using \Cref{coeq for the three diagrams} again.

        It remains to prove that $\mathcal{X}^T$ is finitely cocomplete. We already know that it has finite coproducts, so we only need to show that it has coequalizers. This follows from the fact that it has finite coproducts and reflexive coequalizers, since each coequalizer can be constructed as a reflexive coequalizer from the coproduct, as shown below.
        \[
% https://tikzcd.yichuanshen.de/#N4Igdg9gJgpgziAXAbVABwnAlgFyxMJZABgBpiBdUkANwEMAbAVxiRAEEQBfU9TXfIRQAmclVqMWbAELdeIDNjwEiZAIzj6zVog4BqWTz5LBRURupapuw+JhQA5vCKgAZgCcIAWyRkQOCCQ1ajgACyxXHCQAWmCJbTZXOTdPH0Q-AKCQ8Mis+OsQBxBqBjoAIxgGAAV+ZSEQdywHUKijEA9vJABmakzEUXydEAAdYfwcOgB9YWKQUorq2tNdLDBsWGT21KQBvp6QMIioxGiBqyHkV1I1SekKTY603cDEfcPc-stJC4dr2-uSuVKjUTCpdI1mq0KFwgA
\begin{tikzcd}
A \arrow[rr, "f", shift left] \arrow[rr, "g"', shift right]                       &  & B                                   \\
A+B \arrow[rr, "{[f,1_B]}", shift left=2] \arrow[rr, "{[g,1_B]}"', shift right=2] &  & B \arrow[ll, "\iota_2" description]
\end{tikzcd}
        \]
    \end{proof}  
\end{subbox}

We now want to study a sufficient condition for a category to admit reflexive coequalizers, so that we can apply the proposition above.

Recall that a pair $(f,g)$ from $A$ to $B$ in a regular category induces a relation with the following construction. Consider the morphism induced by the universal property $\langle f,g \rangle$. Since the category is regular, we can take the factorization $\langle f,g \rangle = \langle r_1,r_2\rangle e$, with $\langle r_1,r_2\rangle$ mono and $e$ regular epi. Then $\langle r_1,r_2\rangle$ is the relation induced.
\[
% https://tikzcd.yichuanshen.de/#N4Igdg9gJgpgziAXAbVABwnAlgFyxMJZABgBpiBdUkANwEMAbAVxiRAEEQBfU9TXfIRQAmclVqMWbAEIAdWXgC28AATTuvEBmx4CRAIyl94+s1aIQAJW7iYUAObwioAGYAnCIqRkQOCElEQAAsYOig2HAB3CBCwhGpTKQtWagY6ACMYBgAFfl0hEDcseyCcDVcPL0RAvyRDYNDwxDAmBgZqHDosBgiunoTJcxB5NLB7BhgVNwB9Qynp4RV5NzoxiZBUjKzcnUE2IpKyngrPbw7-RHrEoZHV8cmXUnsl2RW11i4KLiA
\begin{tikzcd}
    A \arrow[rd, "e"', two heads] \arrow[rr, "{\langle f,g \rangle}"] &                                                   & B\times B \\
                                                                      & R \arrow[ru, "{\langle r_1, r_2 \rangle}"', tail] &          
    \end{tikzcd}
\]

\begin{mainbox}{}
    \begin{prop}\label{relation induced by a reflexive pair is reflexive}
        The internal relation induced by a reflexive pair $(f,g)$ in a regular category is reflexive.
    \end{prop}
\end{mainbox}

\begin{subbox}{}
    \begin{proof}
        Consider the pair $(f,g)$ from $A$ to $B$ and the section $s$ with $fs=1_B$ and $gs=1_B$.
\[
% https://tikzcd.yichuanshen.de/#N4Igdg9gJgpgziAXAbVABwnAlgFyxMJZABgBpiBdUkANwEMAbAVxiRAEEQBfU9TXfIRQBGclVqMWbAELdxMKAHN4RUADMAThAC2SMiBwQkokHAAWWNTiQBaAEzV6zVohBruvN1t2J9h49TmltaIDhLObIog1Ax0AEYwDAAK-HgEbBpYimbWPOreAQZGvo6SLqbRILEJyamCbFhg2LByXEA
\begin{tikzcd}
    A \arrow[r, "f", shift left=2] \arrow[r, "g"', shift right=2] & B \arrow[l, "s" description]
    \end{tikzcd}
\]
Let $\langle f,g \rangle = \langle r_1,r_2\rangle e$ be the (regular epi, mono) factorization. Moreover, by the universal property we have that $\langle f,g \rangle s = \langle 1_B,1_B \rangle$.
\[
% https://tikzcd.yichuanshen.de/#N4Igdg9gJgpgziAXAbVABwnAlgFyxMJZAJgBoAGAXVJADcBDAGwFcYkQAhEAX1PU1z5CKMgEZqdJq3YBBHnxAZseAkTIBmCQxZtEnAAQAdQ3gC28fV179lQouVKaa26XqsKlg1SgAsjrVK6nDwSMFAA5vBEoABmAE4QpkiiNDgQSGSSOuzGjPRg4Yww+jGk4UaGcfmFbNYg8YlIDiBpyc6B7Ah1DUmIma2IflmuIMZoWAD6xCA0eQBGMIwACgIqwiBxWOEAFjjysQm9-emI6u3ZemOTovv1h02pJ2cgC2BQSAC06s0uQaITXFm9AWy1Wdj0mx2e2690QzQGQ1e71OPw6en+7gOjUQKRaT3OIxiMxA80WK1s3g2W12tx6bTxSCGv3Y4VpsPhJ0ySM+PgAnASgrlqkV9BjSBiKlUCkViaTQRT1pCadxKNwgA
\begin{tikzcd}
    &  & B \arrow[d, "s"] \arrow[llddd, "1_B"', bend right] \arrow[rrddd, "1_B", bend left] &  &   \\
    &  & A \arrow[dd, "{\langle f,g \rangle}"] \arrow[lldd, "f"] \arrow[rrdd, "g"']                                                                  &  &   \\
    &  &                                                                                                                                             &  &   \\
  B &  & B \times B \arrow[rr, "\pi_2"'] \arrow[ll, "\pi_1"] \arrow[from=uuu, "{\langle 1_B,1_B \rangle}"', bend right=49, crossing over]                                                                                         &  & B
  \end{tikzcd}
\]
Thus $\langle r_1,r_2\rangle es = \langle 1_B,1_B \rangle$. The morphism $es$ makes the relation reflexive.
\[
% https://tikzcd.yichuanshen.de/#N4Igdg9gJgpgziAXAbVABwnAlgFyxMJZARgBoAmAXVJADcBDAGwFcYkQAlEAX1PU1z5CKAAyli1Ok1bsAQjz4gM2PASLlxkhizaIQ83vxVCiZEVum79PSTCgBzeEVAAzAE4QAtkgDMNHBBIYlI67E6GIO5evv6BiGQhMnrEAPryNIz0AEYwjAAKAqrCIG5Y9gAWOAquHt6IfiABSBqJVqkGilF1wU3xNNpJJSnE1ZG1QbHN-ZbsbinkIBnZuQXGanqlFVXclNxAA
\begin{tikzcd}
    & B \arrow[dd, "es"] \arrow[ld, "1_B"'] \arrow[rd, "1_B"] &   \\
  B &                                                         & B \\
    & R \arrow[lu, "r_1"] \arrow[ru, "r_2"']                  &  
  \end{tikzcd}
\]
    \end{proof}
\end{subbox}

\begin{mainbox}{}
    \begin{prop}\label{exact and Maltsev cat has reflexive coeq}
        Let $\mathcal{A}$ be an exact Mal'tsev category. Then $\mathcal{A}$ has reflexive coequalizers.
    \end{prop}
\end{mainbox}

\begin{subbox}{}
    \begin{proof}
        Assume $\mathcal{A}$ is exact and Mal'tsev. We want to prove that $\mathcal{A}$ has reflexive coequalizer. Consider a reflexive pair $(f,g)$ from $A$ to $B$. By \Cref{relation induced by a reflexive pair is reflexive} the internal relation induced $\langle r_1,r_2 \rangle$ is reflexive. Since $\mathcal{A}$ is Malt'sev, then each internal reflexive relation is an equivalence relation. Hence $\langle r_1,r_2 \rangle$ is an equivalence relation. $\mathcal{A}$ is also exact, then, by \Cref{equiv rel is kernel pair of the coeq} the equivalence relation $(r_1,r_2)$ has a coequalizer $k$.
        
        We want to prove that $k$ is also the coequalizer of $(f,g)$. First we note that $kf=kr_1e=kr_2e=kg$. It remains to prove the universal property. Let $h:B \to C$ be a morphism such that $hf=hg$. Then $hr_1e=hr_2e$. Since $e$ is a regular epi, hence an epi, we have that $hr_1=hr_2$. Thus there exists a unique morphism $t$ such that $th=k$.
        \[
    % https://tikzcd.yichuanshen.de/#N4Igdg9gJgpgziAXAbVABwnAlgFyxMJZABgBpiBdUkANwEMAbAVxiRAEEQBfU9TXfIRQBGUgCYqtRizYAhbrxAZseAkTITq9Zq0QgASgr4rBRMeMnaZegDo2AtnRwALAE73gAYQgwAjlwAKVwB9URCxAEojJX5VIWRzYUtpXRBPbkkYKABzeCJQADNXCHskURAcCCQAZi0UtgBrEGoGOgAjGAYABVjTPVcsbOccaKKSpHMKqsRyuGcsApGZlvbOnpM1fsHh5qkdNnDR4tLEMimJuv29VmpnGDooNhwAdwg7h4QeQuOkM8qy6gdMCPRDVM5zBZLAC05SsqQKR3Gp2o-2WICBILB1AhiwBe2sIGyu1aHW6vU2IAGQxGXxAYxOk1Rs3muMQMMuBJCwkRJ3KqIALC0sGBUnAIAwsI8OalnDykILzqChSK2FA6HMsrs4Ww7DAAB5YOA4OAAAgAhCaaRQuEA
    \begin{tikzcd}
        A \arrow[dd, "e", two heads] \arrow[rdd, "f", bend left, shift left] \arrow[rdd, "g"', bend left, shift right] &                                   &                                    \\
                                                                                                                       &                                   & C \arrow[d, "\exists ! t", dashed] \\
        R \arrow[r, "r_2"', shift right] \arrow[r, "r_1", shift left]                                                  & B \arrow[r, "k"'] \arrow[ru, "h"] & {\mathrm{Coeq}(r_1,r_2)}          
        \end{tikzcd}
        \]
    \end{proof}    
\end{subbox}

\section{The notion of ideally exact category}

We are now ready to prove the theorem, which will lead to the definition of the ideally exact categories.

\begin{mainbox}{}
    \begin{theo}\label{Ideally exact categories}.
        Let $\mathcal{A}$ be a category. The following are equivalent:
        \begin{enumerate}
            \item $\mathcal{A}$ is exact and protomodular, has finite coproducts and the the morphism $!:0 \to 1$ is a regular epimorphism.
            \item $\mathcal{A}$ is exact, has finite coproducts and there exists $\mathcal{X}$ semiabelian category with a monadic functor $U: \mathcal{A} \to \mathcal{X}$.
            \[
% https://tikzcd.yichuanshen.de/#N4Igdg9gJgpgziAXAbVABwnAlgFyxMJZABgBpiBdUkANwEMAbAVxiRAB12BbOnACwDGjYAEEAviDGl0mXPkIoyAJiq1GLNpx78hDYAA0JUmdjwEiARnKr6zVog7deg4YYB6wAKoAxI6phQAObwRKAAZgBOEFxIZCA4EEhWanZsniDUcHxYYThIALQWxiCR0UnUCbGZ2bkFybYaDt6S0iVRMYhxlYhK1A32jgIEgRkgfDB0UEhgTAwMFXRYDGyQYKzFpR298YmIyVk5eYiFG+3lO0jbB7XHRRRiQA
\begin{tikzcd}
    \mathcal{A} \arrow[dd, "U", shift left] \arrow[r, "\cong"]      & \mathcal{X}^{UF} \arrow[ldd, shift left] \\
                                                                    &                                          \\
    \mathcal{X} \arrow[uu, "F", shift left] \arrow[ruu, shift left] &                                         
    \end{tikzcd}
            \]
            \item There exist $\mathcal{X}$ semiabelian category with a monadic functor $U: \mathcal{A} \to \mathcal{X}$ and $F$ left adjoint to $U$ such that the associated monad $T=UF$ preserves regular epimorphisms and kernel pairs.
        \end{enumerate}
    \end{theo}
\end{mainbox}

\begin{mainbox}{}
    \begin{defi}
        A category $\mathcal{A}$ is said to be \textbf{ideally exact} if it satisfies one of the equivalent conditions of \Cref{Ideally exact categories}.
    \end{defi}
\end{mainbox}

\begin{subbox}{}
    \begin{proof}[\textit{Proof of \Cref{Ideally exact categories}}]\label{The proof}
        $(1\implies 3)$ Suppose $! : 0 \to 1$ is a regular epi. Since $\mathcal{A}$ is exact, the pullback functor $(!)^*: (\mathcal{A} \downarrow 1) \to (\mathcal{A} \downarrow 0)$ is monadic by (citazione). Note that $(\mathcal{A} \downarrow 1) \cong \mathcal{A}$. Now consider $\mathcal{X} = (\mathcal{A} \downarrow 0)$, we want to prove that $\mathcal{X}$ is semiabelian. 
        
        $(\mathcal{A} \downarrow 0)$ is pointed, in fact if $(A,\alpha)$ is an object in $(\mathcal{A} \downarrow 0)$, we have the following initial and terminal objects, which are the same:
        \[
        % https://tikzcd.yichuanshen.de/#N4Igdg9gJgpgziAXAbVABwnAlgFyxMJZABgBpiBdUkANwEMAbAVxiRGJAF9T1Nd9CKAIzkqtRizYBBLjxAZseAkRFCx9Zq0TtZvRQKIBmUdQ2TtM7nv7KUxtaYladV+XyWDkAFhPjNbDk4xGCgAc3giUAAzACcIAFskERAcCCQAJkd-bQAdHMY0AAs6XRBYhKQyFLTEZLNnAEIAfUs5csTEY2qkAFYs8xA8guLS9srqVIzXMcQ+7sQfPwGhJsC2uI6uyYX+5yGGIpKgziA
        \begin{tikzcd}
        0 \arrow[r, "!_A"] \arrow[rd] & A \arrow[d, "\alpha"] &  & A \arrow[r, "\alpha"] \arrow[d, "\alpha"] & 0 \arrow[ld, "1_0"] \\
                                      & 0                     &  & 0                                         &                    
        \end{tikzcd}
        \]

        $\mathcal{A}$ has finite coproducts, so $(\mathcal{A} \downarrow 0)$ has finite coproducts. In fact consider $(A,\alpha)$, $(B,\beta)$ in $(\mathcal{A} \downarrow 0)$ and $A+B$ in $\mathcal{A}$, with the canonical injections $\iota_1 : A \to A+B$ and $\iota_2 : B \to A+B$. Let $(C,\gamma)$ be an object in $(\mathcal{A}\downarrow 0)$ with $f:A \to C$ and $g:B \to C$ morphisms in $(\mathcal{A}\downarrow 0)$ . Providing to $A+B$ the morphism $[\alpha,\beta]$ we have that $(A+B, [\alpha,\beta])$ is the coproduct of $(A,\alpha)$ and $(B,\beta)$ in $(\mathcal{A} \downarrow 0)$, in fact $\gamma [f,g] = [\alpha,\beta]$, so $[f,g]$ is also a morphism in $(\mathcal{A}\downarrow 0)$.
        \[
% https://tikzcd.yichuanshen.de/#N4Igdg9gJgpgziAXAbVABwnAlgFyxMJZAJgBoAGAXVJADcBDAGwFcYkQBBAagCEQBfUuky58hFGWLU6TVuwDCAoSAzY8BImQCM0hizaIQ5JcLVii5CrtkHOJlSPXjkAFis09cw337SYUAHN4IlAAMwAnCABbJEsQHAgkYkEwyJjELRoEpI8bdgAdfID6KKj6EBpGegAjGEYABUdzQ3CsAIALHHsI6KQAZizExDjPW0L8HHoAfS1utKQ3eKGRvMNxiEmp4gqQKtqGpo0Wts653sQBpaRMkFqwKCQAWj6V-XZQnb26xrMjkFaOl0UiAeulFtkMjQ7g8Lq8vCAAmd0pcIWQZG81vkmGh2uVKjVvodxP8TkDlKCFoMcuj4YVapMkbEqYg0aN2MhCtjcaQ6TBJpRGcNmTdoU8XABOXIYkDIUKkAIC-H7H6iP4A06+fhAA
\begin{tikzcd}
    A \arrow[rr, "\iota_1"] \arrow[rrdd, "f"', bend right] \arrow[rrd, "\alpha"'] &  & A+B \arrow[d] \arrow[d, "{[\alpha,\beta]}"] &  & B \arrow[ll, "\iota_2"'] \arrow[lldd, "g", bend left] \arrow[lld, "\beta"] \\
                                                                                  &  & 0                                                                                 &  &                                                                            \\
                                                                                  &  & C \arrow[u, "\gamma"']  \arrow[from=uu, "{[f,g]}"', bend right=49, crossing over, pos=0.54]                                                          &  &                                                                           
    \end{tikzcd}
    \]

    Now we want to prove that $(\mathcal{A} \downarrow 0)$ is exact, from the exactness of $\mathcal{A}$.
    First we prove it is finitely complete. It suffices to show that it has terminal a object and pullbacks. The terminal object is the one just shown above. It also has pullbacks, by \Cref{U creates con lim}, since pullbacks are connected limits and $\mathcal{A}$ has them.

    Consider now an equivalence relation $(R,r_1,r_2)$ in $(\mathcal{A} \downarrow 0)$, endowed with the morphism $\rho$ and $\alpha$ as shown in the diagram below. We want to prove that such equivalence relation is the kernel pair of some arrow. Since $\mathcal{A}$ is exact, it is a kernel pair in $\mathcal{A}$. By \Cref{equiv rel is kernel pair of the coeq} we can suppose that it is the kernel pair of its coequalizer.
    \[
% https://tikzcd.yichuanshen.de/#N4Igdg9gJgpgziAXAbVABwnAlgFyxMJZABgBpiBdUkANwEMAbAVxiRACUQBfU9TXfIRQBGclVqMWbAILdeIDNjwEiAJjHV6zVohAAhOXyWCio4eK1TdxbuJhQA5vCKgAZgCcIAWyRkQOCCRREDgACyxXHCDNSR0QdwB9VRBqBjoAIxgGAAV+ZSF4rAdQqJ43Tx9EPwDokPDIpABaYMs4xOFDEA9vWprEdQltNgBHFJBQmDooNhwAdwgJqYQyrorfaj6AZhih3QAdPfdQwNSMrNzjFV13IpLO7srgrZ2rEAPGNFC6e7X+jcDENsQAwsGA4lA6GF7GNWmwDjAAB5YOA4OAAQgABAdMjhvlwKFwgA
\begin{tikzcd}
    R \arrow[r, "r_2"', shift right] \arrow[r, "r_1", shift left] \arrow[rd, "\rho"'] & A \arrow[r, "q", two heads] \arrow[d, "\alpha"] & B \arrow[ld, "\exists! \beta", dashed] \\
                                                                                      & 0                                               &                                       
    \end{tikzcd}
    \]
    Since $r_1$ and $r_2$ are morphism in $(\mathcal{A} \downarrow 0)$, we have that $\alpha r_1 = \rho = \alpha r_2$. Thus, by the universal property of the coequalizer, there exists a unique morphism $\beta: B \to 0$ such that $\beta q = \alpha$. Providing to $B$ the morphism $\beta$, we have that $(B,\beta)$ is the coequalizer of $(R,r_1,r_2)$ in $(\mathcal{A} \downarrow 0)$.

    We proceed by proving that $(\mathcal{A} \downarrow 0)$ is regular. Both the conditions of the definition of regular category are easely extendable from $\mathcal{A}$ to $(\mathcal{A} \downarrow 0)$. In fact since in $\mathcal{A}$ each morphism admits a (regular epi, mono) factorization, then the same holds in $(\mathcal{A} \downarrow 0)$. Let $f=me$ be such a fatorization through $Q$, with $f$ morphism in $(\mathcal{A} \downarrow 0)$. Endow $Q$ with the morphism $t$ induced by the unviersal property of the coequalizer. Now we want to prove that $e$ and $m$ are morphism in $(\mathcal{A} \downarrow 0)$. Then they will be respectively regular epi and mono in $(\mathcal{A} \downarrow 0)$. Notice first that $e$ is obviously a morphism in $(\mathcal{A} \downarrow 0)$, by the condition of $t$. Moreover, we have that $te=\alpha=\beta f= \beta me$ and since $e$ is a regular epi, hence an epi, it follows that $m$ is a morphism in $(\mathcal{A} \downarrow 0)$.
    \[
    % https://tikzcd.yichuanshen.de/#N4Igdg9gJgpgziAXAbVABwnAlgFyxMJZABgBoBGAXVJADcBDAGwFcYkQBBEAX1PU1z5CKAEwVqdJq3YAhHnxAZseAkXKliEhizaIQARXn9lQtaRFapukMR4SYUAObwioAGYAnCAFskZEDgQSOqSOuxuRiCePn40gUhiIAAWMPRQ7DgA7hApaQg02tJ6bLzuXr6IifGIITj0WIwZ9Y0FVuy+pVHlsQFBiADMrWF6ADojTGhJ9JHRFbV9g6FFIGMARjB1M92VcQs0jFhg1lD0cCnpQ8tjMAAeWHA4cAAEAIRPOHbcQA
\begin{tikzcd}
    & Q \arrow[rd, "m", tail] \arrow[dd, "\exists ! t", pos=0.75, dashed] &                       \\
A \arrow[rr, "f", pos=0.25, crossing over] \arrow[ru, "e", two heads] \arrow[rd, "\alpha"] &                                                           & B \arrow[ld, "\beta"] \\
    & 0                                                         &                      
\end{tikzcd}
    \]

    In order to prove that regular epis are pullback stable in $(\mathcal{A} \downarrow 0)$ we can use the fact that they are in $\mathcal{A}$. In fact, we just endow the regular epi with the morphism induced by universal property of the coequalizer.

    We must now establish that the $(\mathcal{A} \downarrow 0)$ is protomodular. Since the construction of the change-of-base functor involves pullbacks, which are a connected limits, then by \Cref{connected limits computed on the base} the construction is the same in $(\mathcal{A} \downarrow 0)$ and in $\mathcal{A}$. Thus $(\mathcal{A} \downarrow 0)$ inherits protomodularity from $\mathcal{A}$, noting that the definition of iso in $\mathcal{A}$ and $(\mathcal{A} \downarrow 0)$ is the same. Hence, since $\mathcal{X} = (\mathcal{A} \downarrow 0)$ is exact, protomodular, pointed and has binary coproducts, we can conclude that it is semiabelian.

    It remains to prove that the monad $(!)^*(!)_!$ preserves regular epimorphisms and kernel pairs. Consider a regular epi $q:A \to B$ in $(\mathcal{A} \downarrow 0)$, it is the coequalizer of some parallel arrows. Since the coequalizer is a colimit, it is preserved by the left adjoint functor $(!)_!$. So it suffices to prove that, with $\mathcal{A}$ regular, $(!)^*$ preserves regular epis. Consider the following diagram:
    \[% https://tikzcd.yichuanshen.de/#N4Igdg9gJgpgziAXAbVABwnAlgFyxMJZABgBoBGAXVJADcBDAGwFcYkQBBEAX1PU1z5CKAEwVqdJq3YAhHnxAZseAkXKkREhizaIQ5ef2VCiZYlqm6QxADo28AW3gB9cl15HBq0aXM1t0nq29lhOcK5y3BIwUADm8ESgAGYAThAOSGQgOBBI6pI67ACOhiCp6Zk0OUgiHmVpGYj51Yi1CuWNAMxVuYjEdR1IACw9eQMNSN3ZvSMFgSAAFACEAJQAegBUC0UrPJTcQA
    \begin{tikzcd}
    0\times_1A \arrow[d] \arrow[rr, "(!)^*(q)"] &   & 0\times_1B \arrow[d] \\
    A \arrow[rr, "q"] \arrow[rd]                &   & B \arrow[ld]         \\
                                                & 1 &                     
    \end{tikzcd}
    \]
    By \Cref{p^*(f) is a pullback} the square is a pullback. Since $\mathcal{A}$ is regular, and hence so is $(\mathcal{A} \downarrow 0)$, it follows that $(!)^*$ is a regular epi due to the stability of regular epis under pullbacks.

    Finally, we want to establish that the monad $(!)^*(!)_!$ preserves kernel pairs. This is obvius: $(!)_!$ preserves connected limits, such as kernel pairs, by \Cref{p_! preserves connected limits}, while $(!)^*$ preserves kernel pairs since it is a right adjoint. Thus $(!)^*(!)_!$ preserves kernel pairs.

    $(3 \implies 2)$ Suppose there exists $\mathcal{X}$ semiabelian category and a monadic functor $U: \mathcal{A} \to \mathcal{X}$ satisfying the stated conditions. Assume that the associated monad $T$ preserves
    regular epimorphisms and kernel pairs.
    
    First we need to verify that $\mathcal{A}$ is exact. Since $U$ is monadic, we have that $\mathcal{A} \cong \mathcal{X}^T$. Then it suffies to show that $\mathcal{X}^T$ is exact. We know that $\mathcal{X}$ is semiabelian, so it is exact. Then, by \Cref{algebras category is exact if the base category is exact}, $\mathcal{X}^T$ is exact.

    In order to prove that $\mathcal{A}$ has finite coproducts, it is first necessary to establish that $\mathcal{X}^T$ is protomodular. This follows directly from \Cref{algebras category is proto if the base category is proto} as $\mathcal{X}$ is protomodular. Since any protomodular cateogry is Malt'sev, then $\mathcal{X}^T$ is Malt'sev. We already prove that $\mathcal{X}^T$ is also exact, then by \Cref{exact and Maltsev cat has reflexive coeq}, $\mathcal{X}^T$ has reflexive coequalizers. Since $\mathcal{X}$ has finite coproducts and $\mathcal{X}^T$ has reflexive coequalizers, then by \Cref{algebras category has colim if the base category has coprod}, $\mathcal{X}^T$ has finite coproducts. Finally, since $\mathcal{A} \cong \mathcal{X}^T$, it follows that $\mathcal{A}$  has finite coproducts.


    
\end{proof}
\end{subbox}
